\documentclass[12pt]{article}
\usepackage[margin=2cm]{geometry}
\usepackage{color}
\usepackage{amsfonts}
\usepackage{hyperref}
\hypersetup{
    colorlinks=true,
    linkcolor=blue,
    filecolor=magenta,      
    urlcolor=blue,
    pdftitle={Overleaf Example},
    pdfpagemode=FullScreen,
    }
\pagestyle{empty}

\newcommand{\alert}[1]{{\color{red}#1}}


\begin{document}


%%%%%%%%%%%%%%%% info %%%%%%%%%%%%%%%
\fontsize{12pt}{14pt}\selectfont

\begin{flushright}
    \today
\end{flushright}
\begin{flushleft}
Editors
\\
Physical Review B
\\
American Physical Society
\end{flushleft}
\vspace{-5mm}
%%%%%%%%%%%%%%%%%%%%%%%%%%%%%%%%%%%%%

\noindent\rule{\linewidth}{1.2pt}

%%%%%%%%%%%%%%%%% main %%%%%%%%%%%%%%%
\vspace{5mm}
\fontsize{12pt}{18pt}\selectfont
\noindent Dear Editors,

\bigskip

We wish to submit our research article titled ``\textit{Exact Organization of Density Matrices and Entanglement Structure in the Kitaev Spin Liquid}'' for your consideration in Physical Review B. The article is 19 pages long and includes 8 figures. 

In this research, we introduce a formulation that represents the ground-state density matrix of the spin-1/2 Kitaev honeycomb model in terms of spin operators. Due to the 1-form Wilson symmetry of the phase, both the density matrix and the reduced density matrix in the spin representation can be expressed in an organized form, enabling an exact analysis of the entanglement structure from the perspective of spin degrees of freedom.

The formulation proposed in this research naturally manifests the entanglement structure and its physical meaning related to the emergent $\mathbb{Z}_2$ gauge theory by analyzing the reduced density matrix without the replica trick. In contrast to the replica trick performed using fermionic representation, in which the subtleties of cutting the gauge field on edges and the problems of systems consisting of an odd number of Majorana fermions could arise, our approach provides a unified way to overcome these issues and further reveals the interplay between fermion parity and the entanglement structure.

Recent work in Physical Review X [\href{https://doi.org/10.1103/PhysRevX.15.011001}{W.-T. Xu et al., Phys. Rev. X 15, 011001 (2025)}] has highlighted the important role of generalized symmetries in organizing entanglement properties of gauge theories. 
%
In this context, our work provides an exact microscopic realization of how such organizational structures emerge in a projected gapless parton system. 
%
In particular, our results demonstrate that generalized gauge constraints can strongly organize entanglement structure even in the presence of gapless correlated matter degrees of freedom.

We believe this work meets the editorial standards of Physical Review B for \textbf{conceptual novelty}, \textbf{interdisciplinary interest}, and \textbf{technical rigor}, and we appreciate your consideration.


\vfill
\begin{flushleft}
Sincerely,
\bigskip

Chen-Chih Wang, Yi-Ping Huang, and Sungkit Yip \\

\href{mailto:good.wangchanz@gapp.nthu.edu.tw}{good.wangchanz@gapp.nthu.edu.tw}; \href{mailto:yphuang@phys.nthu.edu.tw}{yphuang@phys.nthu.edu.tw}; \href{mailto:yip@phys.sinica.edu.tw}{yip@phys.sinica.edu.tw}
\end{flushleft}


% %%%%%%%%%%%%%%%%%%%%%%%%%%%%%%%%%%%%%%%
% \newpage

% \alert{
% \href{https://journals.aps.org/prb/authors}{PRB Publishing Guidelines}

% To ensure your papers are handled as Joint Submissions, please clearly indicate your intention in the Cover Letter of all the relevant submissions.

% The cover letter is an opportunity to explain why the manuscript is appropriate for the journal. It should include the following:

% \begin{itemize}
%     \item Context of the results
%     \item Summary of key findings
%     \item Relevant submission history at Physical Review, including another paper that is part of a joint submission
%     \item Any recommended or excluded referees
% \end{itemize}
% }




\end{document}



